%% kapa: exam=1 round=33 year=2022 period=1 form=A subject=real_estate subject_ko=부동산학원론 seq=19 number=99
%% source: Q-Net 감정평가사 기출문제 (cst00309) artlSeq=5209997 file=raw/1008/5209997/2022년도 제33회 감정평가사 1차시험 1교시(민법,경제학,부동산학).pdf
%% conversion: kordoc 4.24.51 -> tools/gen1.py (auto); review=proofread-sonnet; flags=none
\begin{question}{99}
이자율과 할인율이 연 6\%로 일정할 때, A, B, C를 크기 순서로 나열한 것은? (단, 주어진 자료에 한하며, 모든 현금흐름은 연말에 발생함)
\begin{framedtext}
○ A: 2차년도부터 6차년도까지 매년 250만원씩 받는 연금의 현재가치 \par
○ B: 2차년도부터 6차년도까지 매년 200만원씩 받는 연금의 6차년도의 미래가치 \par
○ C: 1차년도에 40만원을 받고 매년 전년대비 2\%씩 수령액이 증가하는 성장형 영구연금의 현재가치 \par
○ 연금현가계수(6\%, 5): 4.212 \par
○ 연금현가계수(6\%, 6): 4.917 \par
○ 연금내가계수(6\%, 5): 5.637 \par
○ 연금내가계수(6\%, 6): 6.975
\end{framedtext}
\begin{choices}
  \item A $>$ B $>$ C
  \item A $>$ C $>$ B
  \item B $>$ A $>$ C
  \item B $>$ C $>$ A
  \item C $>$ B $>$ A
\end{choices}
\end{question}
\answer{4}
